\section{Benchmark Methodology}
\label{sec:methodology}

The purpose of the benchmark methodology is to isolate the effect of
polynomial representation and implementation strategy while avoiding
artificial disadvantages for competing representations. All reported
timings are wall-clock measurements of actual Rust implementations.

\subsection{Benchmark platform}

The headline experiments were run on an Apple M1 system. The repository
records the exact compiler, operating-system, processor, memory, Git
revision, and thread-environment metadata associated with the benchmark in
\texttt{paper/benchmark-environment.txt}.

This machine-generated metadata file is committed with the paper so that
the software environment associated with the measurements can be inspected
independently of the prose description.

The benchmark is compiled in Rust release mode. The release profile uses
one code-generation unit, fat link-time optimization, and abort-on-panic
semantics:
\[
\texttt{codegen-units}=1,
\qquad
\texttt{lto}=\texttt{fat},
\qquad
\texttt{panic}=\texttt{abort}.
\]

Where assembly inspection was required during optimization, compilation was
performed for the native target CPU. The performance results reported in
the benchmark tables are obtained from the release benchmark binaries
described by the repository configuration.

\subsection{Field}

All compared polynomial representations use the same Goldilocks field
implementation:
\[
\mathbb{F}_p,
\qquad
p=2^{64}-2^{32}+1.
\]

Using a common field backend is essential. Comparing one basis with a
specialized field implementation against another basis with generic modular
arithmetic would measure field engineering rather than the representation
itself.

The final Boolean-kernel implementation uses the optimized Goldilocks
arithmetic described in \cref{sec:implementation}. Final cross-basis
numbers must therefore be collected after the same arithmetic backend is
active for every representation participating in the comparison.

\subsection{Input sizes}

The principal scaling experiment uses
\[
N\in
\{
2^{10},
2^{12},
2^{14},
2^{16},
2^{18},
2^{20}
\}.
\]

The headline comparison focuses on
\[
N=2^{20}=1{,}048{,}576.
\]

Power-of-two sizes are natural for the Boolean-kernel tree and also permit
a genuine power-of-two roots-of-unity Lagrange domain in the Goldilocks
field.

\subsection{Input generation}

For reproducibility, the monomial coefficient vector used by the main
Criterion benchmark is deterministic. For
\[
0\le i<N,
\]
the implementation constructs
\[
a_i=17i+3
\]
as a canonical field element.

The Boolean-kernel coefficient vector is obtained from the same polynomial
using the monomial-to-kernel conversion implemented in the repository.
That conversion is performed outside the timed evaluation loop.

The evaluation point used by the main benchmark is
\[
x=19.
\]

These deterministic values are chosen for reproducibility rather than for
special algebraic structure. Correctness tests separately exercise the
representation identities and optimized field arithmetic.

\subsection{Measured operation}

The principal operation is:

\begin{quote}
Given a polynomial already stored in a specified representation and an
arbitrary field point \(x\), compute \(P(x)\).
\end{quote}

The timed region therefore includes the work required to evaluate from that
representation but excludes conversion from another representation.

This distinction is intentional. The benchmark asks how expensive
single-point evaluation is once a prover already owns the polynomial in the
representation being measured. End-to-end systems may need conversion, and
those costs must be accounted for separately when evaluating a complete
prover pipeline.

\subsection{Monomial baselines}

Two monomial baselines are retained.

The first is ordinary serial Horner evaluation. It is useful because it is
the standard direct method and exposes the dependency chain inherent in
that formulation.

The second is a blocked parallel evaluator. If the block size is \(B\),
write
\[
P(X)
=
\sum_j X^{jB}P_j(X),
\]
where each block polynomial \(P_j\) has degree less than \(B\).
The values \(P_j(x)\) can be evaluated independently and combined
afterward.

The paper uses this blocked implementation as the fair parallel monomial
baseline. We do not claim a kernel speedup by comparing a parallel kernel
implementation only against an artificially serialized competitor.

For the headline experiments the monomial block size is
\[
B=4096.
\]

\subsection{Boolean-kernel configuration}

The final Boolean-kernel evaluator uses the implementation described in
\cref{sec:implementation} with local subtree size
\[
B=4096.
\]

The block-size study showed no reproducible advantage large enough to
justify replacing this setting with \(8192\), so \(4096\) is retained as
the conservative default.

Rayon provides the CPU worker pool. Unless a benchmark explicitly studies
thread scaling, the default Rayon worker count for the machine is used.
Thread-scaling experiments, when reported, must state
\texttt{RAYON\_NUM\_THREADS} explicitly.

\subsection{Criterion configuration}

The principal Criterion benchmark uses:

\begin{itemize}
    \item sample size: \(50\);
    \item warm-up time: \(5\) seconds;
    \item measurement time: \(10\) seconds.
\end{itemize}

For close block-size comparisons, a longer controlled benchmark was used:

\begin{itemize}
    \item sample size: \(100\);
    \item warm-up time: \(10\) seconds;
    \item measurement time: \(20\) seconds.
\end{itemize}

We report the Criterion central estimate together with the measurement
interval when discussing individual optimization steps. Outliers reported
by Criterion are not silently removed from the benchmark configuration.

Because sub-millisecond timings on a laptop-class processor can move with
thermal state and background activity, differences at approximately the
one-percent level are treated cautiously and are not promoted to
algorithmic claims without controlled reruns.

\subsection{Correctness checks}

Performance measurements are accepted only after correctness tests pass.

The repository checks, among other properties:

\begin{enumerate}
    \item direct Boolean-kernel evaluation agrees with the tree evaluator;
    \item monomial-to-kernel conversion preserves the represented polynomial;
    \item parallel kernel implementations agree with the serial reference;
    \item optimized Goldilocks multiplication agrees with the reference
    multiplication;
    \item fused kernel folds agree with the unfused mathematical expression;
    \item Lagrange evaluation reconstructs the same polynomial value as the
    monomial representation on test instances.
\end{enumerate}

Optimization variants are therefore compared only after establishing that
they compute the same field value.

\subsection{Treatment of setup and precomputation}

Setup is excluded from a timed loop only when it is reusable in the
corresponding representation-level workload.

For example, converting monomial coefficients into Boolean-kernel
coefficients is not charged to the kernel single-point evaluator because
the benchmark assumes that the polynomial is already represented in that
basis. Likewise, construction of a reusable roots-of-unity domain should
not be charged repeatedly to every Lagrange query if the same domain is
reused by the prover.

Conversely, temporary work intrinsic to one evaluation belongs in the
timed path unless it is explicitly identified as reusable.

All exclusions must therefore be interpreted as representation assumptions,
not as claims that conversion or preprocessing is free.

\subsection{Lagrange comparison policy}

The repository contains both serial and parallel arbitrary-point Lagrange
evaluation over a roots-of-unity domain using batch inversion.

However, the numerical Lagrange results are not frozen at this stage of the
paper. The Boolean-kernel arithmetic backend changed substantially during
optimization, so an old Lagrange measurement collected under an earlier
field implementation is not a valid final comparison.

Before a Lagrange number is placed in the final results table, the Lagrange
implementation will be rerun using the same final field backend, machine,
compiler profile, and Criterion methodology.

This prevents the paper from combining timings collected under materially
different arithmetic implementations.

\subsection{Fairness rules}

The experimental comparison follows the following rules:

\begin{enumerate}
    \item the same field implementation is used by all bases in a final
    comparison;
    \item every method receives the same polynomial size \(N\);
    \item all timings used in a direct ratio are collected on the same
    benchmark machine;
    \item serial and parallel algorithms are labeled separately;
    \item legitimate parallelism is allowed for competing bases;
    \item representation conversion is reported separately from evaluation;
    \item reusable preprocessing is excluded only when the corresponding
    workflow can genuinely reuse it;
    \item all optimized implementations must agree with a correctness
    reference before their timings are accepted;
    \item results within the observed noise floor are not treated as
    meaningful improvements.
\end{enumerate}

These rules are designed to make the benchmark a comparison of useful
algorithms rather than of deliberately favorable implementation choices.
